% !TeX root = ../main.tex
\section{结论}

本文围绕高速公路换电站网络中的预约与随机用户，建立换电路径与充电功率的滚动时域优化方法。候选网络根据 SOC 和道路条件描述车辆可达的换电安排，日前路径提供在线调整的基准。优化模型将路径选择、充电、电池状态及用户服务条件联系起来，以不同周期更新换电路径和充电功率，并按实际服务、充电、路径改变和预约失败记录运营收益。

强化学习用于估计预测域的终端价值，将终端库存与预测域外未完成服务共同纳入评价。模型保留电池 SOC 分布、等待信息及预约后续换电需求，并通过分段线性网络表达其价值。价值函数由完整运营轨迹中的实际回报训练，网络参数在每轮优化中固定。数值实验按联合优化消融、RL 终端价值消融和敏感性分析组织，具体结果与定量结论将在实验完成后填写。

本文以每轮 ETA 预测描述到站窗口，以瞬时换电和给定充电效率表达站内运行。状态编码和价值函数具有有限维度，其效果需要结合实际运营指标和求解时间评价。后续研究可进一步结合交通过程、换电作业时间及电池特性，分析更丰富的运行条件下路径与充电的协调方法。


